\clearpage
\SourceLeaf{321}
\GaussVSourceUnnumbered
\GaussVSourceGathering{40*}
\pdfbookmark[1]{Erdmagnetismus und Magnetometer}{erdmagnetismus-und-magnetometer}
\begin{center}
{\large ERDMAGNETISMUS UND MAGNETOMETER.\par}
\bigskip
\rule{.70\linewidth}{.4pt}\\[.5\baselineskip]
{\small Jahrbuch für 1836 herausgegeben von \textsc{Schumacher}. Stuttgart und Tübingen 1836.}\\[.5\baselineskip]
\rule{.70\linewidth}{.4pt}
\end{center}

Zwei grosse Naturkräfte sind auf der Erde allerorten und in jedem Augenblick
gegenwärtig: die Schwere und die erdmagnetische Kraft.

Die Wirkungen der Schwerkraft sehen wir auf jedem unsrer Schritte uns
begegnen. Die Wirkungen der erdmagnetischen Kraft fallen nicht von selbst in
die Augen, sondern wollen gesucht sein: Jahrtausende vergingen, ohne dass man
nur die Existenz dieser Kraft wusste. Von der erstern Kraft werden alle Verhältnisse
des physischen Lebens durchdrungen, von der andern unmittelbar wenig
oder gar nicht berührt.

Beide Kräfte haben das gemein, dass sie Bewegungen in bestimmten Richtungen
hervorzubringen streben, und dass die Grösse dieser Bewegungen bestimmten
Gesetzen unterworfen ist: aber welche Verschiedenheit, wenn man die
Äusserungen beider Kräfte näher betrachtet!

Zuerst in Beziehung auf die Gegenstände der Kräfte. Der Schwere unterworfen
sind alle materiellen Dinge, vielleicht. und auch nur vielleicht, einige
wenige Stoffe ausgenommen, die man Imponderabilien nennt, und hypothetisch
annimmt, weil wir mit ihrer Annahme eine Unermesslichkeit von Erscheinungen
erklären, und ohne sie nicht erklären können: unter Erklären versteht aber
der Naturforscher nichts anderes, als das Zurückführen auf möglichst wenige und
möglichst einfache Grundgesetze, über die er nicht weiter hinaus kann, sondern
\clearpage
